\clearpage\SourceLeaf{168}\ArithmeticHead{83}{R}{IVb. ALLGEMEINE ELLIPTISCHE FUNKTIONEN.}
\HistoricalFoot{}{11\textsuperscript{\(\ast\)}}
\indent Dagegen erfordert die Bestimmung des Faktors in der Gl. (33) darüber
hinausgehend noch die Kenntnis von Differentialrelationen zwischen den \(\vartheta_{gh}\)
und ihren Nullwerten. Bei der Gleichung (33) braucht man die Relation
\begingroup\makeatletter\tagsleft@true\let\veqno\leqno\makeatother
\begin{equation*}
\vartheta_{11}'(0|x)=\pi p(x)q(x)r(x),\tag{37}
\end{equation*}
\endgroup
\noindent während \textsc{Gauss} (Werke X, 1, S. 310, vergl. \textsc{Geppert} a. a. O.\textsuperscript{83)}, S. 112) die mit
(37) wesentlich äquivalente Gleichung
\begingroup\makeatletter\tagsleft@true\let\veqno\leqno\makeatother
\begin{equation*}
4x\left(p\frac{dr}{dx}-r\frac{dp}{dx}\right)=prq^4\tag{38}
\end{equation*}
\endgroup
\noindent benutzt, in bezug auf die er an jener Stelle nur sagt: »wovon jedoch der Beweis tiefer liegt«\footnote{Siehe die Herleitung der Relation (37) bei \textsc{Jacobi}, \textsc{Jacobis} Werke I, S. 515—517, Art. 4; in bezug auf beide Gln. (37), (38) vergl. \textsc{Weber}, a. a. O.\textsuperscript{102)}, § 20.}. Dieser Beweis findet sich in einer ganz im Jahre 1801
geschriebenen Scheda (Af), und zwar mit Hilfe der Theorie des agM. (Werke
X, 1, S. 209, siehe besonders das »Theorema« auf S. 212), was darum besonders
bemerkenswert ist, weil dadurch die Berechtigung, den Inhalt der aus späterer
Zeit stammenden Aufzeichnung des Handbuchs 16, Bb (Werke X, 1, S. 308
bis 310) hier in Betracht ziehen zu dürfen, gestützt wird.

Soweit hat also \textsc{Gauss} denselben Weg eingeschlagen, den auch \textsc{Jacobi} in
seinen Vorlesungen\footnote{1835/36 ausgearbeitet von \textsc{Rosenhain}, und 1839/40 ausgearbeitet von \textsc{Borchardt}, vergl. die letztere in \textsc{Jacobis} Werken I, S. 497ff.} gewählt hat, um den Schwierigkeiten auszuweichen, die
die für eine direkte Untersuchung des elliptischen Integrals notwendige Integration im komplexen Gebiete darbietet, d. h. er ging von den Thetafunktionen aus. Aber \textsc{Gauss} verfährt im Jahre 1800 sozusagen \textls[100]{moderner} als
\textsc{Jacobi} 1835, indem \textsc{Gauss} das elegante, nach \textsc{Hermite} benannte Verfahren
benutzt, während \textsc{Jacobi} die Transformationsformeln zweiten Grades der Thetafunktionen in etwas umständlicher Weise mit Hilfe der Darstellungen einer
Zahl als Summe von vier Quadraten herleitet\footnote{Diese von \textsc{Jacobi} angewandte Methode findet sich bei \textsc{Gauss} ebenfalls, aber erst später, in dem Handbuch 19, Be (begonnen Mai 1809), S. 147, siehe Werke III, S. 384.}.

Zum Abschluss dieses Teils unseres Berichts haben wir noch den letzten
Satz der Tagebuchnotiz 108. zu besprechen, der wie folgt lautet: »Ferner
haben wir in diesen Tagen die Prinzipien entdeckt, nach denen man die
arithmetisch-geometrische Folge so interpolieren kann, dass man die in der\par
\clearpage\SourceLeaf{169}\ArithmeticHead{84}{L}{L. SCHLESINGER, ÜBER GAUSS’ FUNKTIONENTHEORETISCHE ARBEITEN.}
\noindent gegebenen Folge zu einem beliebigen rationalen Index gehörigen Glieder
durch algebraische Gleichungen zu bestimmen imstande ist.«

So wie der Übergang von einem Gliede \(a_n,b_n\) des Algorithmus (28\,a) des
agM. zum nächstfolgenden \(a_{n+1},b_{n+1}\) dem Übergang von \(x^{2^n}\) zu \(x^{2^{n+1}}\), also
der Multiplikation der Periode \(\varpi'\) mit 2 entspricht, so wird, bei der Interpolation von Gliedern mit einem beliebigen rationalen Index, eine Periode
durch eine ganze Zahl zu teilen sein. \textsc{Gauss} hat also die Gleichungen für
die \textls[100]{Teilung der Perioden} entdeckt. Auch hierüber enthält die Scheda Ac
keinerlei Aufzeichnung; \textsc{Gauss} scheint diese Scheda überhaupt gegen Ende
Mai 1800 abgeschlossen zu haben\footnote{Der Vollständigkeit wegen erwähnen wir noch, dass die letzten Seiten der Scheda Ac eine Aufzeichnung mit der Überschrift »Motus solidi a nullis viribus sollicitati« (Bewegung eines von keinerlei Kräften angetriebenen Körpers) enthalten, in der das Problem bis zu dem dabei auftretenden elliptischen Integral geführt wird. Die Scheda schliesst mit den Worten (Werke X, 1, S. 206); »Die lemniskatischen Funktionen haben wir am 8. Januar 1797 zu betrachten begonnen«. Das bezieht sich auf die Tagebuchnotiz 51. (siehe oben S. 30, 31), in der von dem speziellen lemniskatischen Integral die Rede ist; hier meint \textsc{Gauss} aber offenbar die ganze Theorie der allgemeinen elliptischen Funktionen, die durch die jener speziellen eröffnet worden ist.}. Im selben Jahre hat er dann ein
Handbuch (15, Ba) begonnen, das mit der Werke III, S. 361—374 abgedruckten
Abhandlung über das agM. eröffnet wird; auf diese Abhandlung kommen wir
weiter unten zurück. Zwischen dem Abschluss der Scheda Ac und dem Beginn dieses Handbuchs hat \textsc{Gauss} wohl anderweitige Notizen gemacht, die
aber zum grössten Teil verloren zu sein scheinen. Wir waren bemüht, die
durch diesen Verlust entstandene Lücke durch Bezugnahme auf spätere Aufzeichnungen auszufüllen. Ein ähnliches Verfahren werden wir auch in den
folgenden Erörterungen, die sich auf die Tagebuchaufzeichnungen 109.—111.
beziehen, zu befolgen genötigt sein.
